\begin{mistake}{2026-04-12}{02}

\begin{problemcontent}
设 \(n\) 阶方阵 \(A,B\) 满足 \(A^2-AB=E\)，求 \(r(AB-BA+2A)\)。

\end{problemcontent}

\begin{solutioncontent}
由已知式提取左公因子，得
\[
A(A-B)=E.
\]

由于方阵存在右逆必可逆，故 \(A\) 可逆，且
\[
A-B=A^{-1}.
\]

右乘 \(A\)，得
\[
(A-B)A=EA=E.
\]

展开可得
\[
A^2-BA=E.
\]

而已知又给出
\[
A^2-AB=E.
\]

两式相减，得
\[
AB=BA.
\]

于是
\[
AB-BA+2A=2A.
\]

又因 \(A\) 可逆，所以 \(r(A)=n\)，从而
\[
r(AB-BA+2A)=r(2A)=r(A)=n.
\]

\end{solutioncontent}

\mistakeknowledge{
\begin{itemize}
  \item \textbf{核心公式/定理}：若 \(n\) 阶方阵 \(A\) 满足 \(AX=E\) 或 \(XA=E\)，则 \(A\) 可逆。可逆矩阵满秩，即 \(r(A)=n\)。
  \item \textbf{易错点}：不能一开始就默认 \(AB=BA\)。本题中交换性是由 \(A(A-B)=E\) 推出 \(A\) 可逆后，再严格推得的。
  \item \textbf{关联知识}：矩阵秩在左乘或右乘可逆矩阵时保持不变；数乘非零常数也不改变秩，即 \(r(kA)=r(A)\, (k\neq 0)\)。
\end{itemize}}

\end{mistake}
